\answer{ex:21:1}{(a)~$T=\tau_{\text{rec}}\ln\frac{1}{1-x^*}$：$1.84$~s と $3.68$~s。(b)~$3.36$ から $4.24$~s。感度は $\dd T/\dd x^*=\tau_{\text{rec}}/(1-x^*)=80$~s で、$x^*\to1$ では一斉発火は不規則になる。}
\answer{ex:21:2}{(a)~$8.6$~W、\eqref{eq:energy}と同じ。(b)~$205$~Mbit/s、$0.2$~W。培養（$10^{-4}$~W）の $2\cdot10^3$ 倍。}
\answer{ex:21:3}{$x_{\text{ss}}=1/(1+Ur_b\tau_{\text{rec}})<x^*$ となるのは、$x^*=0.9$ では $r_b>(1/x^*-1)/(U\tau_{\text{rec}})=0.28$~Hz、$x^*=0.99$ では $0.025$~Hz のとき。シナプスの強さは $x_{\text{ss}}$ まで、つまり $10\%$ と $1\%$ 下がる。}
\answer{ex:21:4}{$u(t-k)$ は正規直交で、$m_k=\|P_Vu(t-k)\|^2$。ここで $P_V$ は出力が張る $N$ 次元部分空間への射影である。ベッセルの不等式から $\sum_k m_k\le N$。}
\answer{ex:21:5}{$\tanh$ ありで $\mathrm{MC}\approx9$--$11$、線形リザバーで $15$--$18$（シード3個）。どちらも $\le30$。非線形性の代償は記憶である。}
\answer{ex:21:6}{(a)~$n\le102$。(b)~$5.6$ シグマ。}
\answer{ex:21:7}{未解決の課題。鍵となる対照：読み出しを凍結し、刺激なしの休止後に改善が残るかを確かめる。残らなければ、変わったのは重みではなく状態である。}
